% TOP_PROBLEM: {"id": 3024, "title": "Kirby Problem 5.17", "category_id": 7, "category": {"id": 7, "name": "topology", "display_name": "Topology", "description": "Properties preserved under continuous deformations.", "slug": "topology", "order_index": 7, "created_at": "2026-07-31T15:26:25.671Z"}, "status": "open"}
% FIRST_POSTED: null
% ENTRY_KIND: statement_only
% Imported from: batch14.tex; UnsolvedMath ID 3024
\documentclass[11pt]{article}
\usepackage[margin=25mm]{geometry}
\usepackage{fontspec}
\setmainfont{Latin Modern Roman}
\usepackage{amsmath,amssymb,mathtools}
\usepackage{unicode-math}
\setmathfont{Latin Modern Math}
\usepackage{xurl,hyperref}
\hypersetup{colorlinks=true,urlcolor=blue}
\setcounter{secnumdepth}{0}
\setlength{\parindent}{0pt}
\setlength{\parskip}{5pt}
\setlength{\emergencystretch}{3em}
\newcommand{\sref}[2]{\hyperlink{source:#1:#2}{[S#2]}}
\newcommand{\cataloguescope}[1]{\paragraph{Recorded target.}#1}
\begin{document}
% UnsolvedMath ID 3024; source catalogue code KP-5.17
\setcounter{section}{3023}
\section{Kirby Problem 5.17}\label{3024}
{\small\sffamily UnsolvedMath ID 3024 · Topology}
\subsection{Definitions and mathematical statement}

\paragraph{Shared notation.}$\mathbb N=\{1,2,\ldots\}$, and $\mathbb Z,\mathbb Q,\mathbb R,\mathbb C$ denote the integers, rationals, reals and complex numbers. Any other conventions are specified locally.


Let $(W,\omega)$ be an exact symplectic manifold. A Liouville vector field $V$ satisfies $\mathcal L_V\omega=\omega$, equivalently $d(\iota_V\omega)=\omega$. A Weinstein structure $(\omega,V,\phi)$ additionally has a Morse Lyapunov function $\phi$: $V$ is gradient-like for $\phi$ and $d\phi(V)>0$ away from its critical points, with the usual local gradient condition at them. On a compact domain $V$ is outward transverse to the boundary; on an open Weinstein manifold the exhaustion and completeness conditions are imposed. A Weinstein homotopy is a one-parameter family of such structures allowing isolated ordinary birth/death critical-point degeneracies. In the fixed-form question below $\omega$ is kept fixed and the family remains convex at the boundary (or at infinity).
\paragraph{Statement recorded in the supplied source.}
Given any two Weinstein structures $(W,\omega,V_1,\phi_1)$ and $(W,\omega,V_2,\phi_2)$ on the same symplectic manifold, does there exist a Weinstein homotopy joining them? No flexibility, common Stein complex structure, or simple-connectivity hypothesis is assumed. \sref{3024}{2}
\subsection{Short English statement}
Can any two Weinstein structures for one fixed symplectic form be joined through Weinstein structures?
\subsection{Sources}
\begin{description}
\item[\hypertarget{source:3024:1}{S1}] ULAM.AI and the UnsolvedMath Contributors, \emph{UnsolvedMath: A Curated Collection of Open Mathematics Problems}, record 3024 (KP-5.17). CC BY 4.0. \url{https://huggingface.co/datasets/ulamai/UnsolvedMath}.
\item[\hypertarget{source:3024:2}{S2}] R. İnanç Baykur, Robion C. Kirby and Daniel Ruberman (eds.), \emph{K3: A New Problem List in Low-Dimensional Topology}, Mathematical Surveys and Monographs 295, American Mathematical Society (2026), Problem 5.17, author preliminary version. \url{https://math.berkeley.edu/sites/default/files/surv-295-ruberman-watermarked-author-pdf.pdf}.
\end{description}
\label{end:3024}
\end{document}
